\section{Twist saturation and an open orbit closure}\label{sec:closure}
We complete the proof while allowing the orbit closure to be singular.
The torus saturation, Baire passage, and ambient submersion have
antecedents in \cite[Section on closed-surface dynamics]{Source027}.
Here all of them occur inside the closure of \emph{one} integral orbit.

Let $Y$ be the closure of the image of the $A$-torus in
\eqref{eq:torus-in-closure}. It is irreducible. Choose an irreducible
component $Z_0$ of $Z$ containing $Y$. For each $\gamma\in\F$, some
positive power $\tau_\gamma^{n_\gamma}$ of its Dehn twist preserves
$Z_0$, because each mapping class permutes the finitely many irreducible
components of $Z$.

\begin{lemma}\label{lem:component-nonres}
In a neighborhood of $x$ in $Z_0$, points whose monodromy on every
$\gamma\in\F$ generates its regular centralizer torus are dense in the
complex analytic topology.
\end{lemma}
\begin{proof}
Intersect $Z_0$ with the ambient smooth irreducible-character locus and
with the open where every $\gamma$-monodromy is regular semisimple.
This is a nonempty irreducible open $Z_0'$ containing $x$.
For one curve $\gamma$, pull back the eigenvalue-ordering cover
$T^{\mathrm{reg}}\to(T/W)^{\mathrm{reg}}$ to $Z_0'$. It is finite
\'etale. Each irreducible component of this cover dominates $Z_0'$:
flatness implies that its generic point lies over the generic point of
the integral base. Finiteness makes its image closed, so each component
in fact surjects onto $Z_0'$.

For any nontrivial character of $T$, its being one is a closed condition
on the cover. It cannot hold identically on any component, because every
component has a point over $x$, and every ordering of the eigenvalues at
$x$ is nonresonant. Its zero locus in each component is therefore proper;
under a finite morphism its image is closed and still proper, by dimension.
There are finitely many covering components and countably many characters.
Repeating this for the finite list of curves gives a countable collection
of proper closed subsets of $Z_0'$.

Their complement is dense analytically. For clarity, apply the Baire
theorem on small compact balls in the smooth dense locus of the
irreducible variety $Z_0'$, where each proper algebraic subset has empty
interior; then take closure to include singular points. This also shows
density in every relative analytic neighborhood of $x$.
\end{proof}

\begin{lemma}\label{lem:torus-orbit}
At each point $y$ in the dense set of Lemma~\ref{lem:component-nonres},
the centralizer-torus twist orbit about every $\gamma\in\F$ lies in $Z_0$.
\end{lemma}
\begin{proof}
Cut the surface along the simple curve $\gamma$. For a nonseparating
curve the surface group is an HNN extension of the cut-surface group.
Holding the cut representation fixed, one may multiply a transverse
gluing matrix by any element of the centralizer of $\rho_y(\gamma)$.
This is an algebraic torus action on the gluing parameter and a regular
map from that torus to the character variety. The Dehn twist is the
parameter $\rho_y(\gamma)$ or its inverse, according to orientation.
All curves in our list are nonseparating; the analogous separating
construction would conjugate one cut component.

The orbit under $\tau_\gamma^{n_\gamma}$ lies in $Z_0$. Powers of
$\rho_y(\gamma)^{n_\gamma}$ are still dense in its connected torus:
a character trivial on these powers would give a finite-order value on
$\rho_y(\gamma)$, whose suitable multiple would contradict its full
cyclic density. The inverse image of the closed set $Z_0$ under the torus
map is closed and contains these dense powers, so it is the whole torus.
\end{proof}

We next specify the complex Hamiltonian convention. On an ambient
neighborhood $U$ of $x$ contained in the smooth irreducible-character
locus $D/G$, let $\omega$ be the holomorphic symplectic form from the
complex trace pairing. For a holomorphic invariant function $h_\gamma$
let its holomorphic Hamiltonian field be defined by
\[
 \omega(X_{h_\gamma},\,\cdot\,)=\dd h_\gamma.
\]
The real two-form $\Omega=\Re\omega$ is nondegenerate: if
$\Re\omega(v,w)=0$ for all real tangent vectors $w$, also test $iw$ to
obtain the imaginary part. Under the identification of real tangent
vectors with their $(1,0)$ components, the real Hamiltonian of
$\Re h_\gamma$ has component $X_{h_\gamma}$, and the Hamiltonian of
$\Im h_\gamma$ has component $-iX_{h_\gamma}$.
Indeed these equalities follow directly from
$\Re(-i\dd h)=\Im\dd h$. Choosing $2\Re\omega$ and doubling the real
Hamiltonians gives the same fields.

This is also the variation relation of Goldman
\cite[Lemma 1.11]{Goldman}. His Theorem 4.3, proved in Theorems 4.5 and
4.7 (the HNN formula (4.6) in the nonseparating case), identifies these
real Hamiltonians with centralizer twists. The gradient of an invariant
function belongs to the centralizer Lie algebra. Therefore
$X_{h_\gamma}$ is tangent to the complex centralizer-torus orbit wherever
the monodromy is regular semisimple. This passage uses both real and
imaginary parts and takes place entirely in the ambient smooth locus;
it does not assume smoothness of $Z_0$.

\begin{lemma}\label{lem:ideal}
Every selected holomorphic Hamiltonian field has local complex-time flows
that preserve $Z_0\cap U$, including through its singular points.
\end{lemma}
\begin{proof}
Let $X$ be one such field and let $f$ be a holomorphic function on a small
open subset of $U$ vanishing on $Z_0$. At a point $y$ from
Lemma~\ref{lem:component-nonres}, Lemma~\ref{lem:torus-orbit} and the
Hamiltonian-twist identification give $(Xf)(y)=0$. By analytic continuity
and the density proved there, $Xf$ vanishes on all of $Z_0$ in that open
subset.

The radical ideal sheaf of the reduced analytic subspace $Z_0\cap U$ is
coherent. Near any point choose finite generators $f_1,\ldots,f_s$.
The preceding vanishing says that $Xf_j$ lies in that ideal sheaf; after
shrinking the neighborhood write
\[
 Xf_j=\sum_{k=1}^s a_{jk}f_k
\]
with holomorphic coefficients. Along an integral curve of $X$, the vector
$(f_1,\ldots,f_s)$ satisfies a homogeneous linear differential equation.
With zero initial values its unique solution is zero. Holomorphic ODE
existence and uniqueness, or the corresponding real ODE in each complex
time direction, shows that the local flow preserves the subspace even
when the starting point is singular on $Z_0$.
\end{proof}

\begin{proof}[Proof of Theorem~\ref{thm:main} for $r\geq3$]
At $x$, the chosen invariant differentials span the ambient cotangent
space by Lemma~\ref{lem:choose-x}. Nondegeneracy of $\omega$ makes their
Hamiltonian vectors span the ambient tangent space. Select a basis of
$d=(2g-2)(r^2-1)$ such vectors and compose their local holomorphic flows,
starting at $x$, with complex time parameters $(t_1,\ldots,t_d)$ near zero.
All intermediate points lie in $Z_0$ by Lemma~\ref{lem:ideal}. The
derivative of this map at the origin has the selected basis as its
columns, and is an isomorphism onto $T_x\X$. The holomorphic inverse
function theorem puts a nonempty ambient analytic open subset of $\X$
inside $Z_0$. Since $\X$ is irreducible by the exact source input in
Section~\ref{subsec:geometry}, such an open is Zariski dense. Thus
$Z_0=Z=\X$. The integral seed \eqref{eq:seed} proves the assertion.
The already credited rank-two case completes all $g,r$ in the statement.
\end{proof}
